% SPDX-License-Identifier: CC-BY-4.0
% Copyright (c) 2026 Martin Schr\"oder <info@swedishembedded.com>

\section{Results}
\label{sec:results}

\StatusResults

\subsection{Correctness}
\label{sec:res-correctness}

\input{table-correctness}

\subsection{Learning-rate selection}
\label{sec:res-lr}

\input{table-lr}

Selection picked $10^{-3}$ for SpectralLM and the C=96 pilot and
$3\times10^{-3}$ for every other model (\cref{tab:lr}). SpectralLM's 300-step
losses were not monotone in the rate ($3\times10^{-3}$ was worst of the
three), which one seed per rate cannot separate from seed variation; the rule
was applied as registered.

\subsection{Language modeling at matched width}
\label{sec:res-primary}

\input{table-primary}

\Cref{tab:primary} answers the primary question. Over three seeds SpectralLM
reaches \BpbS{} bits per byte, the parameter-matched SwiGLU control \BpbQpar{}
and the arithmetic-matched SwiGLU control \BpbQstd{}. Against the
parameter-matched control the difference is \DiffSQpar{} bits per byte (Welch
95\% interval \DiffSQparCI{}, $p=\DiffSQparP$ on \DiffSQparDF{} degrees of
freedom). By the registered rule, a mean difference of at most 0.02, H1 is
\HOneVerdict{}; it also holds under the stricter reading that the interval's
upper end lies below the margin: at equal width, equal attention and the same
parameter count, a feed-forward sublayer built from per-mode filters around a
signal-domain gated product predicts held-out text at least as well as a dense
SwiGLU, and its seed range \RangeRelation{}. Three seeds per model do not
establish that it is better: the interval reaches zero.

H2 holds by a wide margin. The SwiGLU of hidden width 1024, with \ParamsQstd{}
non-embedding parameters against \ParamsS{} and nearly the same analytic
arithmetic per token (\MflopQstd{} against \MflopS{} MFLOP with the transforms
counted as FFTs), is better by \DiffSQstd{} bits per byte
(\DiffSQstdCI{}). At this scale the spectral mixer buys parameter efficiency,
not arithmetic efficiency: with the same budget of operations a dense
feed-forward network spends them on four and a half times as many parameters
and uses them better.

The spectral model also costs more to run in this implementation. It trains
at \TpsS{} tokens per second against \TpsQpar{} and \TpsQstd{} for the
controls on the same card, and peaks at \MemS{} MiB of device memory against
\MemQpar{} and \MemQstd{}. Both follow from the mixer's signal-domain
activations: every token carries $C\cdot D=12{,}288$ signal samples per block in
each of three forward buffers, and three more in the backward pass, where the
dense control holds 1024 hidden units; \cref{sec:res-systems} prices the
transforms that produce them.

The pilots, one seed each, bound this result. The spectral mixer's quality
is not monotone in its channel count at the selected rates: $C=8$ scores
\BpbSEight{}, $C=24$ \BpbS{} and $C=96$ \BpbSNinetySix{}, the last at
\MemSNinetySix{} MiB of device memory against \MemS{} for $C=24$, so adding
channels beyond 24 bought nothing measurable here while multiplying the
signal-domain activations. A single seed per pilot cannot say whether $C=96$
is worse than $C=24$ or merely no better. The wide control
$\mathcal{Q}(512,2048)$, with \ParamsQwide{} non-embedding parameters, scores
\BpbQwide{}, the best of the trained models.

\subsection{Post hoc extension}
\label{sec:res-extension}

The registered comparison established non-inferiority but not superiority, and
SpectralLM's rate came from a single-seed selection that was not monotone in
the rate. After the registered analysis we therefore trained three more
seeds of SpectralLM and of its parameter-matched control at their selected
rates, and one SpectralLM run at $3\times10^{-3}$, the rate its selection run
ranked last (the table below). This analysis was not registered.

\input{table-extension}

Over \ExtN{} seeds each, SpectralLM scores \ExtBpbS{} and the control
\ExtBpbQ{}, a difference of \ExtDiff{} bits per byte (Welch 95\% interval
\ExtCI{}, $p=\ExtP$); the interval \ExtVerdict{}. At $3\times10^{-3}$
SpectralLM scores \AltBpb{}, \AltDiff{} against its mean at the selected rate.

\subsection{Comparison with Qwen3-0.6B}
\label{sec:res-qwen}

\input{table-qwen}

Zero-shot, Qwen3-0.6B scores \BpbQwen{} bits per byte on the validation
documents, worse than every model trained on the corpus, the
parameter-matched control included. After continued training on
\QwenFtTokens{} tokens of training stories, an eighth of the trained models'
budget, during which its training loss fell from \QwenFtLossFrom{} to
\QwenFtLossTo{} nats per token, it scores \BpbQwenFt{}, better than every
trained model by a wide margin. We read the zero-shot result as a domain gap:
TinyStories is a narrow synthetic distribution that a small model fitted to it
predicts better than a general-purpose model that has never been adapted to
it, and a few million tokens of adaptation are enough for the larger model's
capacity to dominate. The comparison places the trained models rather than
competing with the pretrained one: SpectralLM has a few million parameters,
two orders of magnitude fewer than Qwen3-0.6B, and has seen 33.6M tokens of one
corpus.

\subsection{Spectral structure of Qwen3-0.6B hidden states}
\label{sec:res-compress}

\input{table-compress}

\Cref{tab:compress} tests the premise behind compressing a pretrained
model's hidden state into a few Fourier coefficients. Along its coordinate
axis, Qwen3-0.6B's residual stream has no spectral structure: the lowest
$m$ Fourier modes keep about the fraction $m/1024$ of the centred energy that
any subspace blind to the coordinates' statistics keeps, and replacing the
stream by that projection raises the next-token loss by several nats at every
width, about as much as the random projection does. As registered, H3's
energy clause fails narrowly: the Fourier basis keeps between 0.013 and 0.021
more energy than the one random draw at every width, while both stay within
0.025 of $m/1024$. Its information clause holds at $m\le128$, where the
Fourier projection costs more loss than the random one, and fails at
$m\ge256$, where it costs less. Its last clause holds by a wide margin: the
top 128 principal components keep almost all of the energy and raise the
loss by less than 0.15 nats.

The table also shows why energy is the wrong criterion. The first 64
coordinates keep 99.5\% of the centred energy, because a few coordinates
carry most of it, yet that projection raises the loss by more than three
nats; what predicts the next token is spread across directions that carry
little energy. Both results agree with \cref{prop:absorption}: a model whose
maps on the residual stream are dense has no reason to order its coordinates
by frequency, so a fixed Fourier truncation of a pretrained model is no
better than a random one, and any spectral representation has to be learned
from the start rather than imposed afterwards.

\subsection{Transform cost}
\label{sec:res-systems}

\input{table-systems}

Both transform paths compute the same model (\cref{tab:correctness}). Where
the FFT path is faster, it shortens the whole training step by a factor of
\FftGainRange{}; the dense basis product is faster at \DenseWins{}, so H4 is
\HFourVerdict{} as registered. By the analytic count a radix-2 transform of
length $D=512$ does about a tenth of the arithmetic of the dense product with
$m=256$ retained modes ($5D\log_2 D$ against $2Dm$ per row), yet the step
shortens by at most the factor above: the transforms are only part of a step,
and the FFT kernels, one workgroup per row with a barrier per radix-2 stage,
use the card far less efficiently than the tiled matrix product they replace.
The fast transform's arithmetic advantage is real but, with these kernels on
this hardware, mostly unrealized.
